\section{结论}
\label{sec:conclusion}

本文针对捕获前自由漂浮空间机械臂，提出了“漏斗预设性能—周期延迟零化”两段式预设时间跟踪方法，并给出了与之配套的扰动观测与切换时序。主要结论如下。

\begin{enumerate}
  \item \textbf{结构层面。}PPC 与 PDF 在语义上正交（前者约束路径、后者约束时间），但结构上不可直接叠加——时变非线性的误差变换会破坏单周期零化所依赖的线性周期时滞结构。据此按时间分段而非按结构叠加，并把交接点取在观测器的预设收敛时刻，使两种机制的激活时刻精确对齐。

  \item \textbf{双重预设。}结论为“预设时间 \(T_o\) 进带 + 再经预设时间 \(2h\) 归零”，总预设时间 \(T_{\mathrm{pre}}=T_o+2h\)，与初始误差和扰动幅值无关（定理~\ref{thm:main}）。进带时刻来自性能函数 \eqref{eq:rho}~的构造，该函数与 PTDO 调节函数同构且在 \(T_o\) 处 C$^{4}$ 光滑，切换不引入由 \(\rho^{(2)}\) 引起的力矩跳变。

  \item \textbf{包络对观测峰值免疫。}漏斗界中 \(\vartheta_i<1\) 只要求观测残差\emph{有限}而与其数值\emph{无关}（定理~\ref{thm:funnel}）。与纯 PD 阶段 I 的输入—状态界对照（推论~\ref{cor:pd_compare}）可见，PPC 的作用不是加速收敛（收敛时间由 PDF 决定），而是把观测峰值从破坏精度降级为消耗进带余量。

  \item \textbf{显式的结构瓶颈。}指出“预设包络在全程成立”并非无条件结论，给出可离线计算的判据 \(M_hB_o\le\min_i\rho_{i\infty}\)（命题~\ref{prop:stage2}），它把人工延迟 \(h\)、Gramian 增益规模与误差带宽度显式耦合，是两段式结构中唯一真正的性能牺牲点。

  \item \textbf{数值验证。}在七自由度浮基机械臂（\(20\,\mathrm{N\cdot m}\) 限幅、\(T_o=5\,\mathrm{s}\)、\(h=1\,\mathrm{s}\)）上，预设性能阶段的包络违反时长为 \(0\)（同包络下纯基准反馈为 \(6.2\,\mathrm{s\cdot joint}\)、归一化比峰值 \(10.6\)），进带深度 \(0.8\%\!-\!5.9\%\)，无扰动终端误差 \(1.2\times10^{-8}\,\mathrm{rad}\)；PDF 驱动阶段的最大越界比 \(1.24\)，加宽误差带后可降至 \(0.69\)，与命题~\ref{prop:stage2}~一致。观测时间扫描表明 \(T_c/T_o\in[0.8,1]\) 为最优区间（阶段 I 饱和占空比由 \(14\%\) 降到 \(0.08\%\) 以下、稳态残差稳定在 \(10^{-4}\,\mathrm{rad}\)），而 \(T_c/T_o\le0.3\) 时抖振使稳态残差恶化三个数量级，且障碍项在饱和下反而带来失稳风险。数值实验还给出两条工程约束：性能包络的初始宽度必须与观测器瞬态相容（第~\ref{sec:sim_infeasible}~节），否则力矩饱和会使障碍项失去意义。

  \item \textbf{边界与降级。}在连续时间、精确模型、精确补偿与无饱和下得到精确归零；在离散实现与数值误差下退化为 \(\Gamma_h\bar\Delta_r/(1-\varepsilon_\Delta)\) 的显式有界形式（命题~\ref{prop:approx}）。同时明确记录三类不可行命题：无扰动观测器时归零不成立且残余带可能大于误差带；\(T_c\le T_o\) 与抖振之间存在定量三角约束 \eqref{eq:tc_tradeoff}；Gramian 病态会使 PDF 退化为等效 PD 并在 \(\varepsilon_\Delta\ge1\) 时失去有界性保证。
\end{enumerate}

后续工作包括：\textup{(1)} 把 \(T_c\) 扫描扩展到参数摄动与初始误差变化的蒙特卡洛统计，并给出 \(T_c/T_o\) 的自整定规则；\textup{(2)} 将性能函数与增益调度结合，使 \(\rho_i\) 或 \(K_0\) 依赖构型指标，从而把本文方法与广义雅可比退化、基座反作用强度等指标联系起来；\textup{(3)} 研究时变基线 \(A_c(t)\) 下的 \(\Phi\)-加权零化增益，以消除切换跳变并容纳时变增益调度；\textup{(4)} 考虑测量噪声下的漏斗实现（噪声会使 \(\lvert\xi\rvert\) 在真值附近抖动，需要给出保包络的噪声容许条件）；\textup{(5)} 在气浮台或硬件在环平台上开展实验验证。
